\documentclass{article}
\usepackage{enumerate}
\usepackage[margin=1.4in]{geometry}
\title{Teaching Statement}
\author{Joel Villatoro}
\date{}

\usepackage{blindtext}
\usepackage{hyperref}
\begin{document}
\vspace{-0.5in}
\maketitle

I consider teaching and mentorship one of my primary and most important responsibilities as an academic mathematician. 
In many ways, it is also the most rewarding and impactful part of my job. 
The core of my teaching philosophy can be described by my two main focuses, student engagement and accessibility. 
By student engagement, I mean getting students to be involved in the course beyond a lecture/homework format. Accessibility means remaining available and approachable to my students and mentees.

Like many mathematicians, my experience teaching began was a teaching assistant leading discussion sections at the University of Illinois at Urbana-Champaign. The main approach there was active learning, which meant that the students were split into small groups and given a worksheet. My role was as a facilitator for discussions and provide students with hints and suggestions if they get stuck. In this time I realized the importance of learning to identify struggling or distracted students and to take proactive action rather than to simply wait for questions.

Unsurprisingly, it is for the larger courses that maintaining student engagement during lectures can be a challenge. 
I have employed a few strategies to break up the tedium of lectures and improve student engagement. My favorite and generally most successful method is the creation of interactive visualizations of mathematical concepts. For example, while teaching differential equations, I created an interactive phase diagram visualizer [link: \url{https://www.desmos.com/calculator/tj3vn0rcqj}] so that students could develop intuition for integral curves and eigenvectors. Another example is python notebook that illustrates a topological proof of the fundamental theorem of calculus [link:\url{https://colab.research.google.com/drive/1A898TEOGSHCNcYKM0BWP8YmsE2PLFz9M?usp=sharing}] which I created for another instructor. 
I am also frequently on the lookout for online resources (such as videos or other multimedia) that can compliment the lectures and worksheets. Especially for undergraduate level topics, there is a vast quantity of high quality visualizations and instructional resources.

Although these strategies help with student engagement, it is my opinion that the discussion sections of large lecture courses are where the most learning actually takes place. At both the University of Illinois and Washington University, active learning is the preferred format. At KU Leuven, I was asked to do a more traditional format of brief lecture followed by solo work on premade worksheets. My conclusion from these experiences is that the active learning approach is much more successful. Without the added stimulation of student-student interactions, discussion sections with solo work or lecture emphasis seem to have significantly less impact.

Although availability is an important component of staying accessible to my students, I believe that it also important to be proactive in getting one-on-one interactions with students especially when they show signs that they are struggling.
After each exam, I always reach out to students who appear to be under performing and invite them to arrange a meeting so we can discuss the course and what they can do to improve.
Especially during one-on-one interactions, I try to remain cognizant of and sensitive to the diverse backgrounds of my students and mentees.
This means taking advantage of common ground I may have with other students as well as respecting the cultural boundaries and personal experiences of my students.

At UIUC, I also gained experience as a mentor for undergraduates through the Illinois Geometry Lab. The Illinois Geometry Lab is a program run by the mathematics department which pairs small groups of undergraduates with a faculty member and graduate student. The faculty member provides the undergraduates with a project, usually a visualization of some mathematical concept, while the graduate student provides regular support while the undergraduates work on the project. The main lesson I took away from this was how rewarding it is to dive into a new topic and learn together with undergraduates. I also learned a lot about programming and numerical methods. Despite the fact that these are not closely related to my research area, I continue to use and develop these skills to enhance my teaching.

Later, at KU Leuven in Belgium I played a variety of teaching roles. In the European system teaching assignments are less standardized and more irregular for postdocs. While there, I taught a graduate course in symplectic geometry in the Spring of 2020. Like many others, I was forced to find technological solutions to the cancellation of classes. This presented a significant challenge but I learned a lot about organizing online lectures and video recordings. At KU Leuven I also worked as a supervisor for a Bachelor project. This involved advising a group of undergraduate mathematics majors while they studied chosen topic in mathematics. In addition to these activities I occasionally worked as a grader for student presentations and occasionally did discussion sections for general education level mathematics classes.

At Washington University in St. Louis, I have taught as the primary instructor for a junior level course in multivariable calculus and a standard course on differential equations. Due to the restrictions placed on me by the NSF grant funding me, I was not permitted to take a larger teaching workload. The multivariable calculus course was for mathematics and physics majors so the class size was relatively small while The differential equations course was much larger (about 150 students between two sections). For the differential equations course, the other section was taught by another professor and we worked together to split the class coordination workload.
\end{document}
